\chapter{Introduzione e requisiti}

PlantPulse è un progetto personale IoT dedicato al monitoraggio e alla gestione delle condizioni ambientali di una pianta. L'intero sistema viene eseguito localmente su un singolo computer. Il dispositivo fisico è, in questa versione, sostituito da un simulatore scritto in C++.

\section{Obiettivi}
\begin{itemize}
    \item simulare l'acquisizione di parametri ambientali;
    \item trasferire le misurazioni al backend tramite HTTP e JSON;
    \item conservare misurazioni, impostazioni e comandi in un database relazionale;
    \item consultare dati correnti e storici tramite una dashboard web;
    \item inviare comandi al dispositivo simulato e registrarne l'esito;
    \item containerizzare i servizi e automatizzare controlli e pubblicazione delle immagini.
\end{itemize}

\section{Requisiti funzionali}
Il sistema gestisce umidità del terreno, luminosità, temperatura dell'aria e umidità dell'aria. Offre quattro aree applicative: Dashboard, Statistics, Controls e Maintenance. Il dispositivo simulato ha identificativo \texttt{plantpulse\_001}.

\section{Perimetro}
Il progetto è una dimostrazione didattica e personale. Non comprende, allo stato attuale, sensori fisici, autenticazione utente, hosting pubblico o deployment su un server remoto.
